% Status: drafted 2026-09-22, merged from the former sections 7 and 8 and cut
% twice the same day. Every number here is also in section 5; this section
% reads them. Each limitation states its consequence for a specific claim.
\section{Discussion and limitations}
\label{sec:errors}
% The old sections 7 and 8 are one section now, and both labels point at it so
% that a forward reference from section 1 or 5 still lands somewhere.
\label{sec:limitations}

\subsection{The accept rule, against the literature and against itself}

The post-correction literature converges on one warning: an LLM corrector
degrades already-good OCR~\cite{kanerva2025nofreelunches}, and over-correction
on low-noise input is HIPE-OCRepair's consistent
problem~\cite{ehrmann2026hipeocrepair}. Both describe a corrector that rewrites
when it should abstain. The marker anchor turns that from a quality problem
into a detectable one: the \corrRewritesCaught{} refusals are replies that
changed text outside the markers, which a rule without the anchor would have
accepted silently. But the branch fires on any byte difference, so it bounds
fluent rewriting from above --- a moved space and a rewritten clause land in
the same bucket, and \corrRewritesSubstantial{} of the \corrRewritesCaught{}
moved more than two characters.

\textbf{A hole in the rule.} It is not as airtight as
section~\ref{sec:system} describes. \texttt{suspect\_spans} locates each
suspect word by character offset when the recogniser's raw text and the page's
text align, and otherwise searches for the word --- taking the \emph{first}
occurrence, so two low-confidence occurrences of the same word resolve to the
same index. The line is then marked with two markers around one word, the model
is shown a corrupted line, and the accept rule compares the reply against that
same corrupted construction, and accepts it. The rule cannot catch this,
because the corruption happens before it runs, and the branch is not an edge
case: \corrOnFindFallback{} of the \corrOffered{} offered lines took it. This
is the most serious limitation in the report. The measured acceptance rate is
an upper bound on the rule's clean behaviour, and \textbf{C2 is supported on
under a third of the lines E2 offered}. One caveat on that number: a failed
offset check is necessary for the bug, not sufficient --- it also needs two
low-confidence occurrences of one word on one line --- so
\corrFindFallbackShare{} bounds the exposure rather than estimating the damage.
The fix belongs in the pipeline: search forward from the previous span's end,
and refuse to mark a line whose spans overlap.

\textbf{And E2 measures refusal, not improvement.} Whether the edits it
accepted made the text better needs a reference transcription for the pages it
sampled, and none exists. Its sample is also \corrPages{} of the
\corrReachPagesWithSpan{} pages carrying a suspect span: a sample of the
corpus, not the corpus.

\subsection{Localisation, and structure}

Livre accuracy is \locLivreAccuracyWorkbook{} and chapter accuracy
\locChapterAccuracyWorkbook{} on the same pieces. The pre-registration
anticipated this case --- \qtr{Correct CHAPTER far below correct LIVRE points
at the span, not the search} --- and said in the same breath to tune
\texttt{CHAIN\_GAP\_PENALTY} and \texttt{WINDOW\_SLACK} before concluding it.
That tuning has not been done, so the span reading is the pre-registered
expectation confirmed in direction, not a tested cause. The denominators say
something the rates do not: only \locChapterScoreableWorkbook{} of
\locPiecesWorkbook{} certain pieces carry a decodable printed chapter number,
so chapter accuracy is the weaker of the two figures in evidence as well as in
value.

On family~A's structure, \structSpuriousA{} proposed chapter openings remain
against \structMissingA{} misses --- after a reader made 128 promotions and 217
demotions by hand. That residue is over-proposal, and it is not evidence about
the unaided detector: the hand pass removed an unknown share of the misses, and
bounding it puts unaided recall as low as 77.0\% against the
\structRecallHandCheckedA{} reported.

\subsection{What is not here}

\textbf{Recognition accuracy is not measured.} This report shows that the
pipeline runs, how fast, what its accept rule refuses and how well it locates a
passage. It does not show how well the fine-tune reads, and it does not compare
it against the published models it replaced. The confusion table the project
has asserted for a year --- long~\qtr{ſ} against \qtr{f}, \qtr{u} against
\qtr{n}, \qtr{c} against \qtr{e} --- remains an assertion. What exists towards
it is the released transcription set of section~\ref{sec:artefacts}, and what
that set cannot settle it states on its own face: it is post-edited rather than
blind, single-annotator, one work and one copy, carrying Transkribus's own
segmentation.

\textbf{Validation is not test.} The one recognition figure the project does
have rests on \trainPagesVal{} held-out pages that selected the shipped
checkpoint, split page-wise from a single volume, so the same formes and the
same type wear appear on both sides. The model has also read some of the text
it was scored on: no page of \emph{Amadis de Gaule} is in the training list,
but the \emph{Trésor} excerpts \emph{Amadis}, and the size of that effect is
not estimable without knowing which passages T.1 excerpts. Every figure in this
report is computed on the corpus rather than on that split, which is why none
of them inherits the problem --- and why none of them is a recognition rate.

\textbf{Five of six thresholds are uncalibrated}
(table~\ref{tab:thresholds}), as are the two the pre-registration names as
preconditions for reading the chapter result. \textbf{And the timings cover the
batch runner only}: \ocrSecondsPerPage{}~s per page is this runner on this GPU,
not a claim about the deployed system, whose execution database was not read.
